\section{Nonlinear stability, open sets, and Gaussian data}
\label{sec:nonlinear}

We first construct a decaying graph for fixed similarity parameters.
By varying phase, center, and blowup time, we then show that every datum
in a small physical $H^k$ neighborhood admits parameters placing its
transformed perturbation on that graph.
The graph construction is a discrete Lyapunov--Perron argument,
related to \cite[Proposition~4.8]{BuckEtAl2026}. Selecting symmetry
parameters by a finite-dimensional topological argument has a focusing
antecedent in \cite[Section~6.3]{Li2023}; here the graph also depends
on the expanding-torus scale.

Fix one sufficiently small admissible $a=1/(p-1)$, with $p$ odd, and one
profile $Q$ given by Proposition~\ref{profile:existence}. Fix the integer
$k>8$ and the spaces, cutoff, and spectral coordinates from
Section~\ref{sec:linear}. Every constant in this section may depend on this
fixed profile and on $k$. In particular, no estimate uniform in $p$ is needed.
We use the Euclidean norm on the fourteen spectral coordinates.

\subsection{Uniform nonlinear step maps}

Let $L_s=L e^{s/2}$ and, on the central cell of the expanding torus, set
\[
 Q_s(y)=\chi(y/L_s)Q(y).
\]
Extend this function periodically. The operator
$\partial_s+y\cdot\nabla/2$ differentiates at fixed $y/L_s$, so the
transport derivative of the cutoff is zero. Write
$\mathcal N(z)=|z|^{p-1}z$. The defect of $Q_s$ in the similarity equation is
\begin{equation}\label{nl:residual-formula}
 \mathcal R_s
 =2\nabla\chi(y/L_s)\cdot\nabla Q
   +\Delta\chi(y/L_s)Q
   +\bigl(\chi(y/L_s)-\chi(y/L_s)^p\bigr)\mathcal N(Q),
\end{equation}
where the derivatives act on the displayed functions of $y$.
Every term is supported in an annulus of radius comparable to $L_s$.
The profile symbol estimates imply, for every fixed integer $j\geq0$,
\[
 |\nabla^j\mathcal R_s|\leq C_j L_s^{-2-2a-j},
 \qquad
 \|\mathcal R_s\|_{Y_{L_s}}\leq C L_s^{-2-a}.
\]
Indeed the low Sobolev order $6-a$ contributes the factor
$L_s^{6-(6-a)}=L_s^a$, and the high order contributes a smaller power.
The regularized zero-frequency part has the same bound. On another cell
write $z=y-2\pi L_s n$. Then
$(\partial_s+y\cdot\nabla/2)z=z/2$, so this computation is compatible
with the periodic lift.

Because $p$ is odd, $\mathcal N$ is a polynomial in $z$ and $\overline z$.
The uniform algebra bounds from Section~\ref{sec:linear} and the uniform
bound for $Q_s$ therefore give, with
\[
 \mathcal E_s(v)=\mathcal N(Q_s+v)-\mathcal N(Q_s)
                    -D\mathcal N(Q_s)v,
\]
the estimates
\begin{equation}\label{nl:quadratic}
 \|\mathcal E_s(v)\|\leq C\|v\|^2,
 \qquad
 \|\mathcal E_s(v)-\mathcal E_s(\widetilde v)\|
 \leq C(\|v\|+\|\widetilde v\|)\|v-\widetilde v\|
\end{equation}
on a fixed small ball. All norms in this display are at the current scale.

Choose the step length $M$ in Proposition~\ref{linear:torus} sufficiently
large that also
\begin{equation}\label{nl:step-condition}
 2e^{-(2+a)M/2}<1.
\end{equation}
For a step beginning at scale $L$, put $L'=Le^{M/2}$ and denote its linear
evolution by $\mathcal F_L$.

\begin{lemma}[Nonlinear step estimate]\label{nl:step}
There are $\delta>0$, $L_0<\infty$ and $C_M<\infty$ such that, for
$L\geq L_0$, the perturbation equation $U=Q_s+v$ has a unique solution
through the step whenever $\|v(0)\|_{Y_L}\leq\delta$. Its endpoint is
\[
 v(M)=\mathcal F_Lv(0)+h_L(v(0)),
\]
where
\begin{align}
 \|h_L(0)\|&\leq C_M L^{-2-a},\label{nl:step-zero}\\
 \operatorname{Lip}_{\|v\|\leq d}h_L
   &\leq C_M(d+L^{-2-a})\qquad(0<d\leq\delta).
 \label{nl:step-lip}
\end{align}
During the step,
\[
 \sup_{0\leq s\leq M}\|v(s)\|_{Y_{L_s}}
 \leq C_M\bigl(\|v(0)\|_{Y_L}+L^{-2-a}\bigr).
\]
\end{lemma}

\begin{proof}
Identify the expanding spaces by dilation and use the free Fourier
propagator from Section~\ref{sec:linear}. On a fixed slab its norm is
bounded uniformly in the starting scale, as are the linearized
multiplication operators. On sufficiently short substeps, the integral
equation is a contraction on a small slab-norm ball by
\eqref{nl:quadratic}; the forcing is bounded by
\eqref{nl:residual-formula}. The short-substep length can be chosen
independently of $L$. Iteration and Gronwall's inequality give the displayed
slab bound for sufficiently small $\delta$ and large $L_0$.

At zero initial perturbation the same integral equation gives
$O_M(L^{-2-a})$, proving \eqref{nl:step-zero}. Differences of two
solutions are at most $C_M$ times their initial difference. Subtract the
linearized solution from this difference. Its remaining source is
$\mathcal E_s(v)-\mathcal E_s(\widetilde v)$; use
\eqref{nl:quadratic} and the slab bound. Integration over the fixed slab
gives \eqref{nl:step-lip}. These arguments also prove uniqueness.
\end{proof}

\subsection{A graph of decaying perturbations}

Set $L_n=Le^{nM/2}$ and abbreviate $E_{L_n},\pi_{L_n},P_{L_n}$ by
$E_n,\pi_n,P_n$. The normalized coordinates satisfy $\pi_nE_n=I$.
Let $D_*=e^{M\Lambda}$, where $\Lambda$ has entries $0$, twelve copies
of $1/2$, and $1$. Thus $\|D_*^{-1}\|\leq1$.
Split
\[
 v_n=w_n+E_nu_n,\qquad w_n\in\ker\pi_n,\quad u_n\in\mathbb R^{14}.
\]
Proposition~\ref{linear:torus} and Lemma~\ref{nl:step} give the recursions
\begin{align}
 w_{n+1}&=A_nw_n+B_nu_n+(I-P_{n+1})h_{L_n}(v_n),\label{nl:forward}\\
 u_{n+1}&=D_*u_n+e_n(v_n)+\pi_{n+1}h_{L_n}(v_n),\label{nl:coordinates}
\end{align}
where
\[
 \|A_n\|\leq\tfrac18,\qquad \|B_n\|\leq C_M,\qquad
 \sup_{n\geq0}\|e_n\|\longrightarrow0\quad\text{as }L\to\infty.
\]
Here $e_n=\pi_{n+1}\mathcal F_{L_n}-D_*\pi_n$. All coordinate maps
and projections are uniformly bounded.

\begin{proposition}[Continuous stable graph]\label{nl:graph}
For all sufficiently large $L$ there is a continuous map $G_L$, defined
on a fixed small ball in $\ker\pi_L$, such that the initial perturbation
\[
 v_0=w_0+E_LG_L(w_0)
\]
generates a solution on every similarity-time slab and
\[
 \|v(s)\|_{Y_{L_s}}\longrightarrow0\quad(s\to\infty).
\]
Uniformly in $L$,
\begin{equation}\label{nl:graph-size}
 |G_L(w_0)|\leq C\bigl(\|w_0\|+L^{-2-a}\bigr).
\end{equation}
The coordinate $G_L(w_0)$ is jointly continuous in $L$ and $w_0$ after
identifying periodic functions by dilation.
\end{proposition}

\begin{proof}
For a constant $K\geq1$ to be fixed, use the sequence norm
\[
 \|(w,u)\|_{\mathcal X}
 =\sup_{n\geq0}2^n\bigl(\|w_n\|+K|u_n|\bigr).
\]
Define a map by the right side of \eqref{nl:forward} for $w_{n+1}$,
resetting its zeroth coordinate to the prescribed $w_0$, and by
\begin{equation}\label{nl:backward}
 u_n=-\sum_{j=n}^{\infty}D_*^{n-j-1}
       \bigl(e_j(v_j)+\pi_{j+1}h_{L_j}(v_j)\bigr)
\end{equation}
for the coordinate sequence. The geometric weights make this sum
convergent even in the phase coordinate, where $D_*$ equals one.

For differences of sequences, the linear terms in the forward map cost
at most $2(1/8+C_M/K)$ in the sequence norm. Choose $K$ so that this is
less than $1/2$. The weighted sum in \eqref{nl:backward} costs at most
$2K$ times the weighted supremum of the forcing difference. By
\eqref{nl:step-lip}, the remaining Lipschitz costs tend to zero as the
sequence ball is made small and then $L$ is made large. The same is true
of $e_j$, uniformly in $j$. Consequently the complete map has a
contraction constant strictly below one, uniformly for all sufficiently
large $L$.

At the zero sequence the residual forcing is $O_M(L_n^{-2-a})$.
Condition~\eqref{nl:step-condition} bounds its weighted sequence norm by
$C_M L^{-2-a}$. Thus a fixed small ball is invariant if $w_0$ is small,
and the contraction theorem yields a unique sequence in that ball.
Its norm is at most $C(\|w_0\|+L^{-2-a})$. The fixed point satisfies
\eqref{nl:coordinates}, since consecutive backward sums telescope.
Define $G_L(w_0)=u_0$. This proves \eqref{nl:graph-size} and the decay
at step endpoints. Lemma~\ref{nl:step} gives the decay between endpoints.

We spell out the continuity assertion because the free Schr\"odinger
flow is not continuous in operator norm as its time varies. At every
fixed index, and after dilation to a fixed torus, the free Fourier
formula is strongly continuous in the starting scale and in the input.
The polynomial integral equation preserves this continuity on each
finite slab. Cutoffs and coordinate maps are likewise strongly
continuous: prove this first for smooth functions, then use density and
their locally uniform bounds. The finite matrix normalizing $\pi_L$
stays invertible and varies continuously.

Start the sequence contraction at zero. Each iterate is continuous at
every fixed coordinate: in \eqref{nl:backward}, the summands are
continuous and the geometric tail is uniform. The iterates converge
uniformly in the sequence norm on the allowed parameter set. Their
zeroth coordinate is therefore continuous. This proves the assertion
without differentiating the graph or taking an operator-norm limit of
Schr\"odinger propagators.
\end{proof}

\subsection{Physical parameters fill an open neighborhood}

For a small $T_0>0$, define the smooth periodic reference datum by
\begin{equation}\label{nl:reference}
 u_{\mathrm{ref}}(x)=T_0^{-a}\chi(x)Q(x/\sqrt{T_0})
\end{equation}
in the central cell. Its support is contained strictly inside that cell.
For physical data $u_{\mathrm{in}}$ near this datum, vary
\[
 T=T_0e^\beta,\qquad x_*=\sqrt{T_0}\,\zeta,\qquad \theta\in\mathbb R,
 \qquad d'=(\theta,\zeta,\beta)\in\mathbb R^{14}.
\]
Put $L=T^{-1/2}$. The required starting perturbation is
\begin{equation}\label{nl:starting-perturbation}
 v_0(d';u_{\mathrm{in}})
 =e^{-i\theta}T^a u_{\mathrm{in}}(x_*+\sqrt T\,y)
      -\chi(y/L)Q(y).
\end{equation}

\begin{lemma}[Parameter expansion]\label{nl:parameter-expansion}
For $|d'|\leq\rho$ with $\rho$ fixed and small, uniformly as $T_0\to0$,
\begin{equation}\label{nl:leading-expansion}
 v_0(d';u_{\mathrm{ref}})
 =E_L C d'+r(d'),\qquad
 C d'=(-\theta+b\beta,\zeta,\beta),
\end{equation}
and
\[
 \|r(d')\|_{Y_L}\leq C\bigl(|d'|^2+\sqrt{T_0}|d'|\bigr).
\]
For fixed $T_0$, replacing $u_{\mathrm{ref}}$ by $u_{\mathrm{in}}$
adds at most $C(T_0)\|u_{\mathrm{in}}-u_{\mathrm{ref}}\|_{H^k}$,
uniformly on the parameter ball. All these maps are continuous.
\end{lemma}

\begin{proof}
Substitution of \eqref{nl:reference} into
\eqref{nl:starting-perturbation} gives
\[
 e^{-i\theta}e^{a\beta}
 \chi(\sqrt{T_0}\zeta+\sqrt T\,y)
 Q(\zeta+e^{\beta/2}y)-\chi(\sqrt T\,y)Q(y).
\]
Taylor expansion of the smooth reference profile yields
\[
 \chi(y/L)\bigl(-i\theta Q+\zeta\cdot\nabla Q
                    +\beta(a+y\cdot\nabla/2)Q\bigr).
\]
The profile symbol bounds and their dyadic norm estimates apply to its
first two parameter derivatives: translations preserve the tail bounds,
and the dilation derivative $y\cdot\nabla$ preserves their order.
Hence the Taylor remainder is $O_{Y_L}(|d'|^2)$ uniformly in $L$.
The translated cutoff differs by $O(\sqrt{T_0}|\zeta|)$, with the same
scaled derivative estimates on its annulus; multiplication by the
profile gives the additional stated error. For small parameters the
supports remain strictly inside one physical cell.

The spectral time vector is $(a-ib+y\cdot\nabla/2)Q$.
The missing $ibQ$ contributes $b\beta$ to the phase coordinate, proving
\eqref{nl:leading-expansion}. On the compact parameter ball with $T_0$
fixed, translation and rescaling from physical $H^k$ to $Y_L$ have
uniformly bounded norms. They are strongly continuous, which proves the
remaining assertions. Only the smooth reference has been differentiated;
no extra derivative is required of nearby $H^k$ data.
\end{proof}

\begin{proof}[Proof of Theorem~\ref{thm:blowup}]
Choose $a$ and $k$ as above. We seek a zero of the continuous map
\begin{equation}\label{nl:degree-map}
 \Phi(d')=\pi_Lv_0(d';u_{\mathrm{in}})
       -G_L\bigl((I-P_L)v_0(d';u_{\mathrm{in}})\bigr).
\end{equation}
The matrix $C$ in Lemma~\ref{nl:parameter-expansion} is invertible.
At every current scale, $\pi_LE_L=I$ and $(I-P_L)E_L=0$ exactly.
Consequently the stable argument of $G_L$ contains only the error in
that lemma; varying $L$ introduces no leading stable component.
By \eqref{nl:graph-size}, on $|d'|\leq\rho$ we have
\begin{equation}\label{nl:degree-error}
 |\Phi(d')-Cd'|
 \leq C_1\bigl(\rho^2+\sqrt{T_0}\rho
                   +\varepsilon_{\mathrm{data}}+T_0^{1+a/2}\bigr),
\end{equation}
where
$\varepsilon_{\mathrm{data}}=C(T_0)
\|u_{\mathrm{in}}-u_{\mathrm{ref}}\|_{H^k}$.

Let $c_C>0$ be the smallest singular value of $C$.
First choose $\rho$ small enough for the quadratic term in
\eqref{nl:degree-error} to be less than $c_C\rho/4$ and for the graph
domain requirements to hold. Next choose $T_0$ sufficiently small that
all scales are in the range of Proposition~\ref{nl:graph}, and the two
$T_0$ terms together are less than $c_C\rho/4$.
Finally choose a positive physical $H^k$ radius so that the data term
is less than $c_C\rho/4$. The graph is then defined throughout the
closed parameter ball, and on its boundary
\[
 |\Phi(d')-Cd'|<c_C\rho\leq |Cd'|.
\]
The straight-line homotopy to $C$ has no boundary zero. Brouwer degree
therefore gives a zero of \eqref{nl:degree-map} for every datum in this
open $H^k$ ball.

At such a zero the initial perturbation lies on the graph. The resulting
similarity solution exists for all $s\geq0$ and its perturbation tends
uniformly to zero, by Proposition~\ref{nl:graph} and the uniform
$Y_{L_s}\to L^\infty$ embedding. Undoing the similarity variables gives
the physical solution on $0\leq t<T$:
\[
 u(t,x)=e^{i\theta}(T-t)^{-a}
       \left(\frac{T-t}{T}\right)^{ib}
       \bigl(Q_s(y)+v(s,y)\bigr),
 \quad y=\frac{x-x_*}{\sqrt{T-t}},\quad
 s=\log\frac{T}{T-t}.
\]
On each finite similarity slab, the norms are equivalent after dilation
to the physical $H^k$ norm. The integral equation and polynomial algebra
therefore give $u\in C([0,T);H^k)\cap C^1([0,T);H^{k-2})$.
The usual local uniqueness follows directly from the physical
Schr\"odinger integral equation and its $H^k$ difference estimate.
At the center, $Q_s(0)=Q(0)\ne0$, and hence
\[
 (T-t)^a|u(t,x_*)|\longrightarrow |Q(0)|>0.
\]
Since $H^k(\mathbb T^{12})$ embeds continuously into $C^0$, the solution
cannot extend continuously in $H^k$ through $T$. This proves the theorem.
\end{proof}

\paragraph{Higher-dimensional deterministic examples.}
For the fixed odd power $p$ in Theorem~\ref{thm:blowup}, smooth
finite-time blowup data also exist on every standard product torus
$\mathbb T^d$ with $d\ge12$. To see this, choose the smooth reference
datum \eqref{nl:reference} at the center of the blowup ball just
constructed. Its solution $u$ remains smooth before its blowup time $T$.
Indeed, for each integer $m\ge k$, the Fourier convolution estimate gives
\[
 \big\||v|^{p-1}v\big\|_{H^m(\mathbb T^{12})}
 \le C_{m,p,k}\|v\|_{H^k(\mathbb T^{12})}^{p-1}
                 \|v\|_{H^m(\mathbb T^{12})}.
\]
The frequency weight of order $m$ is bounded by a sum of weights on
single factors, and the remaining factors are controlled by
$\|\widehat v\|_{\ell^1}\lesssim\|v\|_{H^k}$ since $k>6$.
On every $[0,\tau]\subset[0,T)$, the $H^k$ norm of $u$ is bounded.
Duhamel's formula and Gronwall therefore bound each higher Sobolev norm;
local $H^m$ existence, continuation, and $H^k$ uniqueness identify these
solutions with $u$ throughout $[0,\tau]$. The equation then supplies all
time derivatives. This persistence argument takes place entirely in
dimension twelve.

Now set $U(t,x,y)=u(t,x)$ on
$\mathbb T^{12}\times\mathbb T^{d-12}$. The extra Laplacian derivatives
vanish, so $U$ is a smooth solution of \eqref{eq:nls}.
The polynomial nonlinearity is Lipschitz on bounded sets, so the $L^2$
difference estimate and Gronwall give uniqueness among smooth solutions
on compact time intervals.
For every $y$, the center limit from Theorem~\ref{thm:blowup} becomes
$(T-t)^{1/(p-1)}|U(t,x_*,y)|\to c_0>0$.
Thus the same time and pointwise blowup rate persist along the product
fiber. No higher-dimensional stability or Gaussian-probability conclusion
is asserted.

\subsection{The Gaussian counterexample}

\begin{proof}[Proof of Corollary~\ref{cor:gaussian}]
Fix the pair $p,k$ supplied by Theorem~\ref{thm:blowup} and take
$\alpha>k+6$. Independence and the finite Gaussian second moment give
\[
 \mathbb E\|u_0^\omega\|_{H^k}^2
 =c_g\sum_{n\in\mathbb Z^{12}}\langle n\rangle^{2k-2\alpha}<\infty,
\]
up to the fixed Fourier normalization. Thus the series converges as an
$H^k$-valued Gaussian random variable. Its finite Fourier segment has a
strictly positive density on the whole finite-dimensional complex
coefficient space.

For completeness, let $B_{H^k}(f,\varepsilon)$ be any nonempty ball.
Choose a finite Fourier projection $P_N$ such that
$\|(I-P_N)f\|_{H^k}<\varepsilon/4$ and
\[
 \mathbb E\|(I-P_N)u_0^\omega\|_{H^k}^2<\varepsilon^2/32.
\]
Markov's inequality gives positive probability that the random tail
has norm less than $\varepsilon/4$ after increasing $N$ further if
necessary. Independently, the finite segment belongs to the ball of
radius $\varepsilon/2$ about $P_Nf$ with positive probability.
The intersection lies in $B_{H^k}(f,\varepsilon)$, so the law has full
support in $H^k$.

Apply this to the open blowup neighborhood in Theorem~\ref{thm:blowup}.
It has positive probability for every $\alpha>k+6$. Since
$8<k<\alpha-6$, the exponent $m=k$ lies in the interval in the universal assertion of
Corollary~\ref{cor:gaussian}. The local solution fails its required global continuation
with positive probability. The fixed pair also satisfies
$s_{\mathrm{cr}}=6-2/(p-1)>1$. Hence no finite threshold for this pair,
and therefore no universal assertion of the stated form, is possible.
\end{proof}
